\section{Related Work}

\noindent\textbf{Vision-based tactile sensing.}
GelSight-type sensors convert elastomer deformation into images under structured illumination~\cite{johnson2009retrographic,yuan2017gelsight}; GelSlim~\cite{donlon2018gelslim,sipos2024gelslim4} and DIGIT~\cite{lambeta2020digit} adapt the principle to fingertips. Because the measurement depends jointly on gel mechanics, LED layout and camera response, models transfer poorly between sensors; sensor-invariant representations~\cite{gupta2025sitr,zhao2024transferable,higuera2024sparsh,higuera2025sparshx} and deformation-space translation~\cite{zai2025across,yang2025flow} attack this on the representation side. We use a frozen sensor-invariant encoder~\cite{gupta2025sitr} as backbone so that our comparison isolates the data.

\noindent\textbf{Tactile simulation.}
Physically based tactile renderers synthesize GelSight images from mesh contacts~\cite{agarwal2021simulation,wang2022tacto,luu2023simulation}, with example-based~\cite{si2022taxim}, particle-based~\cite{chen2023tacchi}, differentiable~\cite{xu2022efficient} and GPU-scale~\cite{nguyen2024tacex,akinola2025tacsl} variants. Learned renderers instead map simulated depth or contact state to sensor images with an image-to-image network~\cite{isola2017pix2pix,church2022tactile,zhong2025tactgen}. Our two renderers follow the two lines and share meshes, contact sampler and labels, so they differ only in the appearance model.

\noindent\textbf{Unpaired refinement.}
When labeled real pairs are unavailable, synthetic images can be refined toward the real distribution with an adversarial loss on unlabeled real images: SimGAN~\cite{shrivastava2017simgan} adds a self-regularization term to keep the annotation valid, CycleGAN~\cite{zhu2017cyclegan} a cycle constraint, CUT~\cite{park2020cut} a patchwise contrastive constraint, and RCAN~\cite{james2019rcan} translates randomized renders to a canonical one. In tactile sensing, CycleGAN-based translation~\cite{chen2022cyclegan,jing2023uada}, SightGAN's background- and contact-reconstruction terms~\cite{azulay2024sightgan} and depth-conditioned diffusion~\cite{higuera2023braille} have been applied to GelSight-type sensors, evaluated on classification or contact position. These works do not state which real frames the translator saw relative to the test set, nor charge the translator's real data against the downstream budget. Self-training with pseudo labels~\cite{lee2013pseudo,xie2020noisy} is the standard way to use unlabeled data in vision; our mesh-aligned pseudo-labeling (Sec.~\ref{sec:mapl}) applies it to both sides of a real-to-sim-to-real pipeline, the renderer and the model, and replaces each pseudo label by the object mesh at a pose found with a learned residual aligner, as in iterative render-and-compare pose refinement~\cite{li2018deepim,wang2020self6d} but trained from the $\Kreal$ labeled tactile frames alone; our refiner (Sec.~\ref{sec:refiner}) adapts the unpaired family to the finding of this paper---geometry transfers, contact shading does not---by fixing the contact geometry through a mask composite and translating only the shading. Both are evaluated under the same leakage-controlled budget as the paired renderers.

\noindent\textbf{Sim-to-real transfer and its evaluation.}
Domain randomization~\cite{tobin2017domain,peng2018sim} and structured randomization around a calibrated center~\cite{chen2024general} trade fidelity for robustness; sensor-background conditioning~\cite{zhong2025tactgen} and low-fidelity pretraining~\cite{gano2025lowfidelity} report gains from synthetic tactile data. These works measure the benefit of simulation but rarely report how much real data the simulator consumed, or whether the simulator saw the test distribution. Cross-sensor transfer is usually attacked on the representation side~\cite{gupta2025sitr,zhao2024transferable}; we instead ask how much of a new sensor's or session's real data a renderer needs. No prior tactile study, to our knowledge, compares a physically calibrated and a learned renderer on the same meshes, contact sampler and downstream model, with the renderer's real data held disjoint from the test frames. Our contribution is this protocol---renderer and downstream model trained on the same disjoint real subset, tested on held-out objects, a second session and a second sensor---the resulting data-efficiency curves for depth, normals, pose and classification with both renderer families, and mesh-aligned pseudo-labeling, a self-training method with mesh-aligned pseudo labels that the curves motivate and that exceeds the leaky gain at $\Kreal=100$ from unlabeled frames.

\noindent\textbf{Pseudo-labeling and render-and-compare.}
Self-training on a model's own predictions~\cite{lee2013pseudo,xie2020noisy} is the standard way to use unlabeled data, but its labels inherit the model's errors. Object pose estimation refines a pose hypothesis by comparing a render of the model with the image, iteratively and with a learned update~\cite{li2018deepim}, and can self-supervise on unlabeled real images~\cite{wang2020self6d}. SnapTac brings the two together for tactile sensing: the pseudo label of a frame is not a network output but the object mesh at a pose found by a learned residual aligner, and the aligner is trained from the few labeled tactile frames alone, without a differentiable or photorealistic renderer.
